% !TeX root = ../main.tex

\chapter{Results}\label{chapter:results}

\section{Dataset}

We evaluate InsertGS on a curated set of indoor 3DGS reconstructions
distributed as \ac{SPZ} files, including living rooms, restaurants,
libraries, and basements. Inserted objects are textured meshes; the canonical
example used throughout the development of the system is a single
manufacturing chair.

% TODO: replace these placeholders with the real numbers + figures from
% src/cases/*/output/.

\section{Qualitative comparison}

\Cref{fig:qual} (TODO) shows three views of the same room with the chair
inserted using (a) flat shading, (b) shading with a generic studio HDR, and
(c) shading with the InsertGS cubemap extracted from the room itself.
Method~(c) reproduces the room's dominant warm tones on the chair, while
(a) and (b) leave the chair looking pasted on.

\section{Ablations}

\begin{itemize}
  \item \textbf{Without EWA splatting.} Replacing the proper EWA rasterizer
        with naive point splatting in Stage~1 produces visible aliasing at
        cubemap seams (TODO: figure).
  \item \textbf{Without shadow pass.} Removing the shadow-catcher pass
        causes the chair to appear to float; reintroducing it restores
        contact (TODO: figure).
  \item \textbf{Without depth gating.} Foreground Gaussians no longer
        occlude the chair, breaking parallax in animated views (TODO:
        figure).
\end{itemize}

\section{Performance}

End-to-end runtime for a single case on the development workstation
(RTX 4090) breaks down roughly as follows:

\begin{itemize}
  \item Cubemap extraction: $\sim$2~s per insertion point.
  \item Blender PBR rendering: 10--30~s per view at 1080p depending on
        sample count.
  \item Compositing: real time per view once the host 3DGS scene is
        loaded.
\end{itemize}

The \texttt{batch} subcommand parallelizes across cases at the
case-directory level rather than within a case, which matches how the
Blender step parallelizes internally.
